\section{Small cells and fiber mass}\label{sec:smallcells}

We now combine the subset estimate with the overlap identity.  The goal is
to show that cells of size at most $n^{1-c}$ carry only $o(W+1)$ of the
off-diagonal overlap.  The main issue is a cell containing a source but
only a small fraction of its fiber; the paired coordinate will control
this contribution.

First, Equation~\eqref{eq:additive} and the uniform marginals of $\pi$
give $W_2=b_*+2A$.  If
$(q_X+q_Y)/2\ge(1-\epsilon_n)\alpha$, where $\epsilon_n\to0$, then
\[
 \E_q L^*(x,y)
 =W_2+\sum_{x\in P}
   \bigl(q_X(x)+q_Y(x)-2\alpha(x)\bigr)S^*(x).
\]
Since $S^*\ge0$ and $q\le B_0\pi$, it follows that
\begin{equation}\label{eq:q-overlap}
 W_2-o(1)A
 \le \I\bigl(q(C^2)\bigr)
 =\E_q L^*(x,y)
 \le B_0W_2.
\end{equation}
For a current cell occurrence $C$, write $a=\alpha(C)$ and
$p=\mu_e(C)$; the latter depends on the fiber $e$.  Define
\[
 T(\lambda)=\I\left(
   \E_e\1_C(x_e)p\,\1_{\{p<\lambda\}}\right).
\]
This is finite, since its integrand is at most
$q(C^2)=\E_e\1_C(x_e)\mu_e(C)$.

\begin{lemma}\label{lem:small-fiber}
For the prepared graph and its hierarchies, if
$0<\lambda<1/2$, $2/n\le\tau\le1$, and $H\ge1$, then
\begin{equation}\label{eq:small-fiber}
 T(\lambda)\ll
 \left(\frac{\lambda}{\tau^2}+\lambda\right)W_2
 +\tau A+\lambda H+e^{-b_1H},
\end{equation}
where $b_1>0$ is absolute.  The implicit constant depends only on $B_0$
and the absolute constants in Equation~\eqref{eq:grid-error}.
\end{lemma}

\begin{proof}
We separate large cells, complements of giants, and actual small cells.
For a set of size $m\ge2$,
\begin{equation}\label{eq:sc-pair-comparison}
 \frac{\alpha(C)^2}{2}
 \le \pi(C^2)=\frac{m(m-1)}{n(n-1)}
 \le\alpha(C)^2.
\end{equation}
If $a\ge\tau$, then $m\ge2$ and the integrand of $T(\lambda)$ is
at most $\lambda\le2\lambda\pi(C^2)/\tau^2$.  These occurrences
contribute at most $2\lambda W_2/\tau^2$.
If $C$ is the complement occurrence of a giant and $a<\tau$, its
contribution is at most
\[
 q(C^2)\le B_0\pi(C^2)\le B_0\tau a.
\]
The integral of $a$ over the giant-complement occurrences is $A$,
so their contribution is at most $B_0\tau A$.
These two classes retain the multiplicities of the transformed family.

It remains to treat actual nongiant cells of mass less than $\tau$.
Fix a place, a coordinate $j$, and the two entire realized hierarchies.
Let $j'$ be the other coordinate and set $H_v=H$ at archimedean places,
$H_v=0$ otherwise.  For a fiber $e$ with source $x=x_e$, let $U$
be the actual $j$-cell containing $x$ at level $r$, and put
$p=\mu_e(U)$.  We retain only the levels with $U$ nongiant,
$\alpha(U)<\tau$, and $p<\lambda$.
Without grid errors, a fiber point outside the source cell in coordinate
$j$ lies inside the source cell at the complementary level $h_e-r$ in
coordinate $j'$.  We use two nearby levels in the latter coordinate to
measure failures of this complementarity.  Define the cells
\[
 V_-=\Pi_{j'}(h_e-r-H_v)(x),\qquad
 V_+=\Pi_{j'}(h_e-r+H_v)(x),
\]
where $\Pi_j(t)(x)$ denotes the cell of $x$ in $\Pi_j(t)$.
Thus $V_+\subseteq V_-$.  The two errors are
\[
 D^-=\mu_e\bigl(P\setminus(U\cup V_-)\bigr),\qquad
 D^+=\mu_e(U\cap V_+).
\]

The fiber identity makes these errors small after integration.  For a
fixed fiber point $y$, write $d=d_j(x,y)$ and $d'=d_{j'}(x,y)$.
Its contribution to $D^-$ occurs when
$d<r<h_e-H_v-d'$, and its contribution to $D^+$ occurs when
$h_e+H_v-d'\le r\le d$.  Hence
\begin{align*}
 \int D^-\,dr&=\E_{y\sim\mu_e}(h_e-d-d'-H_v)_+,\\
 \int D^+\,dr&=\E_{y\sim\mu_e}(d+d'-h_e-H_v)_+.
\end{align*}
Their sum is the expectation of
$(|\Delta_j(x,y)+\Delta_{j'}(x,y)|-H_v)_+$, since
$i_j(x,y)+i_{j'}(x,y)=h_e$.  At nonarchimedean places it is zero.
At archimedean places Equation~\eqref{eq:grid-error} and the union bound
give
\[
 \Pr\bigl(|\Delta_j+\Delta_{j'}|>u\bigr)
 \le2C_0e^{-b_0u/2}.
\]
Integrating this tail, averaging over fibers, and using
$\sum_{v\mid\infty}\sigma_v=1$, we obtain
\begin{equation}\label{eq:sc-errors}
 \avg{\E_e\int(D^-+D^+)\,dr}\ll e^{-b_1H}.
\end{equation}
No independence of the two errors is required.

We now split the retained levels into three cases.  In each change of
level below the fiber is fixed, so $h_e$ is constant; both realized
hierarchies are held fixed over all levels before any averaging.

\smallskip\noindent
\emph{Case 1: $V_-$ is nongiant.}
Since $\mu_e(V_-)+D^-\ge1-p>1/2$,
\[
 p\le2\lambda\bigl(\mu_e(V_-)+D^-\bigr).
\]
After discarding the other restrictions, change levels to
$r'=h_e-r-H_v$.  Sharing a nongiant cell in the $j'$-hierarchy
contributes to $L_{j'}$, so the main term is bounded by
$2\lambda\E_qL_{j'}$.  Summing over coordinates and places gives
at most $2B_0\lambda W_2$.  The error is bounded by
Equation~\eqref{eq:sc-errors}, since $2\lambda<1$.

\smallskip\noindent
\emph{Case 2: $V_-$ is giant and $V_+$ is not.}
For a fixed source $x$, the levels at which its cell in the
$j'$-hierarchy is giant form an initial interval, by nesting.
The two levels defining $V_-$ and $V_+$ straddle its endpoint, so
the allowed $r$ have length at most $2H_v$.  As $p<\lambda$, this
costs at most $2\lambda H_v$ per fiber.  Summing and averaging gives
at most $4\lambda H$.

\smallskip\noindent
\emph{Case 3: $V_+$ is giant.}
This is the remaining case, since $V_+\subseteq V_-$.  Here
\[
 p=D^++\mu_e(U\setminus V_+).
\]
The first term is covered by Equation~\eqref{eq:sc-errors}.
In the second change levels to $r'=h_e-r+H_v$ for each fiber.
At a fixed new level, every counted pair has
$y\notin G_{j'}(r')$ and shares an actual $j$-cell of size less than
$\tau n$ at some original level.

For a fixed target $y$, let $N_j(y)$ be the union of all actual
$j$-cells containing $y$ whose size is less than $\tau n$, over all
levels.  These cells form a nested chain of subsets of the finite set
$P$, so $|N_j(y)|<\tau n$ (with empty union allowed).
Thus all the possible sources of this target lie in
$N_j(y)\setminus\{y\}$, even though their selected levels depend on
the fiber.  Each edge belongs to exactly one fiber and retains weight
$1/|E|$ after the level change.  The mass at $r'$ is therefore at most
\begin{align*}
 &\sum_{y\notin G_{j'}(r')}
   \sum_{x\in N_j(y)\setminus\{y\}}q(x,y)\\
 &\hspace{1cm}\le
 \frac{B_0}{n(n-1)}
 \sum_{\substack{y\notin G_{j'}(r')\\N_j(y)\ne\varnothing}}
       (|N_j(y)|-1)
 \le B_0\tau\alpha\bigl(P\setminus G_{j'}(r')\bigr).
\end{align*}
In the last step, $|N_j(y)|-1\le\tau(n-1)$ because $\tau\le1$.
Integrating, summing coordinates and averaging gives $B_0\tau A$.

The three cases, the initial split, and Equation~\eqref{eq:sc-errors}
prove Equation~\eqref{eq:small-fiber}.  Threshold ties affect only
endpoints of level intervals.  All rearranged integrands are
nonnegative, so Tonelli's theorem applies.
\end{proof}

We can now discard the small cells in the overlap scale.  Recall
$\lambda_n=1/\log n$.  Taking
$\tau=\lambda_n^{1/3}$ and $H=\lambda_n^{-1/2}$ in
Lemma~\ref{lem:small-fiber} gives
\[
 T(\lambda_n)
 \ll\lambda_n^{1/3}W+\lambda_n^{1/2}
       +e^{-b_1\lambda_n^{-1/2}}
 =o(W+1).
\]
Define
\[
 W_{2,\mathrm{sm}}=
 \I\bigl(\pi(C^2)\1_{\{|C|\le n^{1-c}\}}\bigr).
\]
Split $\I(q(C^2))$ according to whether
$\mu_e(C)<\lambda_n$.  Equation~\eqref{eq:subset}, whose errors are
uniform in every subset $C$, bounds the other part by
\[
 (1+o(1))(W_2-W_{2,\mathrm{sm}})+o(1)W_{2,\mathrm{sm}}.
\]
Adding $T(\lambda_n)=o(W+1)$ and comparing with
Equation~\eqref{eq:q-overlap} yields
\begin{equation}\label{eq:small-overlap}
 W_{2,\mathrm{sm}}=o(W+1).
\end{equation}
This estimate holds before passing to either a bounded or an unbounded
subsequence of $W$.
